
% =============================================================================
% Coxeter, Vinberg and cusp diagrams
% =============================================================================
% Every named diagram is a pic. Its vertices are nodes named by their index, so
%   \pic (A) at (6,0) {k3 cusp=1};
% names the vertices A0, A1, ...; without a pic name they are 0, 1, ....
% Highlight a parabolic subdiagram through the vertex centres:
%   \scoped[on background layer] \draw[parabolic] (A0.center) -- (A8.center);
%
% `root labels=<scheme>` labels every vertex by its index:
%   r      $r_i$  (the default value of a bare `root labels`)
%   ell    $\ell_i$
%   index  $i$
%   prime  $i'$
%   none   no labels (the default)
%
% Published names:
%   k3 cusp=1..6, k3 cusp=2 cover   Coxeter diagrams at the 0-cusps of F_{En,2}
%                                   and their K3 covers
%   sterk=1..5                      Sterk's diagrams I-V of the same five cusps
%   vinberg 18 2 0, vinberg 18 0 0  Vinberg diagrams of (18,2,0) and (18,0,0)
%   k3 F2, k3 F4                    Coxeter diagrams for F_2 and F_4
%   \coxeterHthree \coxeterHfour \coxeterGtwo \coxeterGtwoAffine
%   \coxeterpair[<keys>], \coxetervertex{<mark>}   table legends

\newif\ifdzg@rootlabels
\newif\ifdzg@barbelllabels
\dzg@barbelllabelstrue
\def\dzg@labelscheme{none}
\def\dzg@ellscheme{ell}
\tikzset{
  root labels/.is choice,
  root labels/r/.code={\dzg@rootlabelstrue\def\dzg@labelscheme{r}},
  root labels/ell/.code={\dzg@rootlabelstrue\def\dzg@labelscheme{ell}},
  root labels/index/.code={\dzg@rootlabelstrue\def\dzg@labelscheme{index}},
  root labels/prime/.code={\dzg@rootlabelstrue\def\dzg@labelscheme{prime}},
  root labels/none/.code={\dzg@rootlabelsfalse\def\dzg@labelscheme{none}},
  root labels/.default=r,
}
\def\dzg@label@r#1{$r_{#1}$}
\def\dzg@label@ell#1{$\ell_{#1}$}
\def\dzg@label@index#1{$#1$}
\def\dzg@label@prime#1{$#1'$}
\def\dzg@rootlabel#1{\csname dzg@label@\dzg@labelscheme\endcsname{#1}}

% \dzg@vertex{<index>}{<mark>}{<coordinate>}{<label direction>}
\def\dzg@vertex#1#2#3#4{%
  \ifdzg@rootlabels
    \node[vertex=#2, label={[vertex label]#4:\dzg@rootlabel{#1}}] (#1) at (#3) {};%
  \else
    \node[vertex=#2] (#1) at (#3) {};%
  \fi}

% \dzg@vertices{<index>/<mark>/<coordinate>/<label direction>, ...}
% Brace a coordinate that contains a comma; place a calc coordinate with
% \dzg@vertex, since \foreach does not strip the braces around it.
\def\dzg@vertices#1{%
  \foreach \dzg@i/\dzg@m/\dzg@p/\dzg@a in {#1}{\dzg@vertex{\dzg@i}{\dzg@m}{\dzg@p}{\dzg@a}}}

% \dzg@edges{<Coxeter weight>}{<path through vertex names>}
\def\dzg@edges#1#2{\scoped[on edge layer] \draw[coxeter edge=#1] #2;}

% \dzg@chain{<Coxeter weight>}{<first>}{<last>}: edges k -- k+1, first <= k < last.
% A path \foreach through node names restarts every segment at the first node,
% so each edge is its own path.
\def\dzg@chain#1#2#3{%
  \pgfmathtruncatemacro\dzg@stop{#3-1}%
  \foreach \dzg@k in {#2,...,\dzg@stop}{%
    \pgfmathtruncatemacro\dzg@n{\dzg@k+1}%
    \dzg@edges{#1}{(\dzg@k) -- (\dzg@n)}}}

% \dzg@squareperimeter{<side>}{<edges per side>}{<first index>}{<sw|nw>}
%   {<mark at corners and every second vertex>}{<mark between>}
% Vertices around the square [0,side]^2, numbered from the south-west corner
% counterclockwise (sw) or from the north-west corner clockwise (nw), joined
% in a cycle of simple edges.
\def\dzg@squareperimeter#1#2#3#4#5#6{%
  \def\dzg@L{#1}%
  \pgfmathtruncatemacro\dzg@last{#3+4*#2-1}%
  \foreach \dzg@k in {#3,...,\dzg@last}{%
    \pgfmathtruncatemacro\dzg@s{div(\dzg@k-#3,#2)}%
    \pgfmathtruncatemacro\dzg@t{mod(\dzg@k-#3,#2)}%
    \pgfmathsetmacro\dzg@u{#1*\dzg@t/#2}%
    \pgfmathsetmacro\dzg@v{#1-\dzg@u}%
    \csname dzg@square@#4\endcsname
    \ifodd\dzg@t\def\dzg@m{#6}\else\def\dzg@m{#5}\fi
    \dzg@vertex{\dzg@k}{\dzg@m}{\dzg@p}{\dzg@a}}%
  \dzg@chain{3}{#3}{\dzg@last}%
  \dzg@edges{3}{(\dzg@last) -- (#3)}}
% Position \dzg@p and outward label direction \dzg@a of the vertex at step
% \dzg@t along side \dzg@s.
\def\dzg@square@corner#1#2{\ifnum\dzg@t=0 \def\dzg@a{#1}\else\def\dzg@a{#2}\fi}
\def\dzg@square@sw{%
  \ifcase\dzg@s
    \edef\dzg@p{\dzg@u,0}\dzg@square@corner{225}{270}%
  \or\edef\dzg@p{\dzg@L,\dzg@u}\dzg@square@corner{315}{0}%
  \or\edef\dzg@p{\dzg@v,\dzg@L}\dzg@square@corner{45}{90}%
  \or\edef\dzg@p{0,\dzg@v}\dzg@square@corner{135}{180}%
  \fi}
\def\dzg@square@nw{%
  \ifcase\dzg@s
    \edef\dzg@p{\dzg@u,\dzg@L}\dzg@square@corner{135}{90}%
  \or\edef\dzg@p{\dzg@L,\dzg@v}\dzg@square@corner{45}{0}%
  \or\edef\dzg@p{\dzg@v,0}\dzg@square@corner{315}{270}%
  \or\edef\dzg@p{0,\dzg@u}\dzg@square@corner{225}{180}%
  \fi}

% \dzg@barbell{<corner>/<inner> for the sw, se, ne, nw corners}{<top>}{<bottom>}
%   {<corner edge weight>}{<half length of the barbell>}
% Inside the square [0,4]^2: white inner vertices at (1,1), (3,1), (3,3), (1,3)
% joined to their corners, and a black barbell (infinity edge) at x=2 joined by
% weight-4 edges to the sw and ne inner vertices (top) and to the se and nw
% inner vertices (bottom).
\def\dzg@barbell#1#2#3#4#5{%
  \begingroup
    \ifdzg@barbelllabels\else\dzg@rootlabelsfalse\fi
    \dzg@vertex{#2}{black}{2,2+#5}{0}%
    \dzg@vertex{#3}{black}{2,2-#5}{0}%
  \endgroup
  \dzg@edges{infinity}{(#2) -- (#3)}%
  \foreach \dzg@c/\dzg@i [count=\dzg@k] in {#1}{%
    \ifcase\dzg@k\or\def\dzg@p{1,1}\def\dzg@a{100}\or\def\dzg@p{3,1}\def\dzg@a{80}%
    \or\def\dzg@p{3,3}\def\dzg@a{90}\or\def\dzg@p{1,3}\def\dzg@a{80}\fi
    \ifodd\dzg@k\def\dzg@b{#2}\else\def\dzg@b{#3}\fi
    \dzg@vertex{\dzg@i}{white}{\dzg@p}{\dzg@a}%
    \dzg@edges{#4}{(\dzg@c) -- (\dzg@i)}%
    \dzg@edges{4}{(\dzg@i) -- (\dzg@b)}}}

% \dzg@trianglepoints{<A>}{<B>}{<C>}{<n>}{<m>}{<p>}: coordinates A, B, C and
% the subdivision points ab-0..n of AB, bc-0..m of BC, ca-0..p of CA.
\def\dzg@trianglepoints#1#2#3#4#5#6{%
  \coordinate (A) at (#1);
  \coordinate (B) at (#2);
  \coordinate (C) at (#3);
  \foreach \dzg@k in {0,...,#4}{\coordinate (ab-\dzg@k) at ($(A)!\dzg@k/#4!(B)$);}
  \foreach \dzg@k in {0,...,#5}{\coordinate (bc-\dzg@k) at ($(B)!\dzg@k/#5!(C)$);}
  \foreach \dzg@k in {0,...,#6}{\coordinate (ca-\dzg@k) at ($(C)!\dzg@k/#6!(A)$);}}

% -----------------------------------------------------------------------------
% The five 0-cusps of F_{En,2} (k3 cusp=1..5) and the K3 covers of their
% diagrams (k3 cusp=6 covers cusps 1, 3, 4, 5; k3 cusp=2 cover covers cusp 2).
% White: roots of square -2; black: roots of square -4.
% -----------------------------------------------------------------------------
\tikzset{pics/k3 cusp/.style={/tikz/transform shape, code={%
  \tikzset{transform shape=false}\csname dzg@kcusp@#1\endcsname}}}

\@namedef{dzg@kcusp@1}{%
  \dzg@squareperimeter{4}{2}{0}{sw}{black}{black}%
  \dzg@vertices{8/black/{0.75,0.75}/above left, 9/black/{1.5,1.5}/above left,
    10/black/{2.25,2.25}/above left, 11/black/{3,3}/above left}%
  \dzg@edges{3}{(0) -- (8) (11) -- (4)}%
  \dzg@edges{infinity}{(8) -- (9) -- (10) -- (11)}}

\@namedef{dzg@kcusp@2}{%
  \dzg@trianglepoints{0,0}{4,0}{2,4}{5}{6}{6}%
  \dzg@vertices{0/white/bc-6/above, 1/black/ca-1/left, 2/black/ca-2/left,
    3/black/ca-3/left, 4/black/ca-4/left, 5/black/ca-5/left,
    6/black/ab-0/below, 7/black/ab-1/below, 8/black/ab-2/below}%
  \dzg@vertex{9}{black}{$(A)!0.3!(B)!0.2!(C)$}{below}%
  \dzg@edges{3}{(1) -- (2) -- (3) -- (4) -- (5) -- (6) -- (7) -- (8) (6) -- (9)}%
  \dzg@edges{4}{(0) -- (1)}}

\@namedef{dzg@kcusp@3}{%
  \dzg@vertices{0/white/{0,0}/below, 1/white/{4,4}/above,
    2/black/{0,1}/left, 3/black/{0,2}/left, 4/black/{0,3}/left, 5/black/{0,4}/left,
    6/black/{1,4}/above, 7/black/{2,4}/above, 8/black/{3,4}/above,
    9/black/{1,3}/right}%
  \dzg@vertex{10}{black}{$(0,0)!1/3!(4,4)$}{left}%
  \dzg@vertex{11}{black}{$(0,0)!2/3!(4,4)$}{right}%
  \dzg@edges{3}{(2) -- (3) -- (4) -- (5) -- (6) -- (7) -- (8) (5) -- (9)}%
  \dzg@edges{infinity}{(10) -- (11)}%
  \dzg@edges{4}{(8) -- (1) (2) -- (0) (0) -- (10) (11) -- (1)}}

\@namedef{dzg@kcusp@4}{%
  \dzg@vertices{0/white/{0,2}/left, 1/black/{0,3}/left, 2/black/{0,4}/left,
    3/black/{1,4}/above, 4/black/{2,4}/above, 5/black/{3,4}/above,
    6/black/{4,4}/above, 7/black/{4,3}/right, 8/white/{4,2}/right,
    9/black/{1,3}/below, 10/black/{3,3}/below}%
  \dzg@edges{3}{(1) -- (2) -- (3) -- (4) -- (5) -- (6) -- (7) (2) -- (9) (6) -- (10)}%
  \dzg@edges{4}{(1) -- (0) (7) -- (8)}}

\@namedef{dzg@kcusp@5}{%
  \dzg@squareperimeter{4}{2}{0}{sw}{black}{black}%
  \dzg@barbell{0/8, 2/9, 4/10, 6/11}{13}{12}{4}{0.4}}

\@namedef{dzg@kcusp@6}{%
  \dzg@squareperimeter{4}{4}{0}{sw}{doubled}{white}%
  \dzg@barbell{0/16, 4/17, 8/18, 12/19}{20}{21}{4}{0.4}}

\@namedef{dzg@kcusp@2 cover}{%
  \dzg@trianglepoints{0,0}{4,0}{2,4}{5}{6}{6}%
  \dzg@vertices{16/doubled/ab-0/below left, 17/white/ab-1/below,
    18/doubled/ab-2/below, 2/doubled/ab-3/below, 3/white/ab-4/below,
    4/doubled/ab-5/below right,
    5/white/bc-1/right, 6/doubled/bc-2/right, 7/white/bc-3/right,
    8/doubled/bc-4/right, 9/white/bc-5/right, 10/doubled/bc-6/above,
    11/white/ca-1/left, 12/doubled/ca-2/left, 13/white/ca-3/left,
    14/doubled/ca-4/left, 15/white/ca-5/left}%
  \dzg@vertex{1}{white}{$(A)!0.7!(B)!0.2!(C)$}{above}%
  \dzg@vertex{19}{white}{$(A)!0.3!(B)!0.2!(C)$}{above}%
  \dzg@chain{3}{2}{18}%
  \dzg@edges{3}{(16) -- (19) (1) -- (4)}}

% -----------------------------------------------------------------------------
% Sterk's diagrams I-V (Sterk 1991) of the same five cusps, numbered 1', 2', ...
% -----------------------------------------------------------------------------
\tikzset{pics/sterk/.style={/tikz/transform shape, code={%
  \tikzset{transform shape=false}\csname dzg@sterk@#1\endcsname}}}

\@namedef{dzg@sterk@1}{%
  \dzg@squareperimeter{6}{2}{1}{nw}{black}{black}%
  \dzg@vertices{9/black/{1.5,1.5}/170, 10/black/{2.5,2.5}/170,
    11/black/{3.5,3.5}/170, 12/black/{4.5,4.5}/170}%
  \dzg@edges{3}{(3) -- (12) (7) -- (9)}%
  \dzg@edges{infinity}{(9) -- (10) -- (11) -- (12)}}

\@namedef{dzg@sterk@2}{%
  \dzg@vertices{1/black/{2,0}/270, 2/black/{1,0}/270, 3/black/{0,0}/270,
    4/black/{0,1}/170, 5/black/{0,2}/170, 6/black/{0,3}/170, 7/black/{0,4}/170,
    8/black/{0,5}/170, 9/white/{0,6}/100, 10/black/{1,1}/10}%
  \dzg@edges{3}{(1) -- (2) -- (3) -- (4) -- (5) -- (6) -- (7) -- (8) (3) -- (10)}%
  \dzg@edges{4}{(8) -- (9)}}

\@namedef{dzg@sterk@3}{%
  \dzg@vertices{1/white/{0,0}/180, 2/black/{0,1.5}/180, 3/black/{0,3}/180,
    4/black/{0,4.5}/180, 5/black/{0,6}/135, 6/black/{1.5,6}/90,
    7/black/{3,6}/90, 8/black/{4.5,6}/90, 9/white/{6,6}/90,
    10/black/{1.5,4.5}/10, 11/black/{2,2}/100, 12/black/{4,4}/100}%
  \dzg@edges{3}{(2) -- (3) -- (4) -- (5) -- (6) -- (7) -- (8) (5) -- (10)}%
  \dzg@edges{infinity}{(11) -- (12)}%
  \dzg@edges{4}{(2) -- (1) -- (11) (8) -- (9) -- (12)}}

\@namedef{dzg@sterk@4}{%
  \dzg@vertices{1/white/{0,3}/180, 2/black/{0,4.5}/180, 3/black/{0,6}/135,
    4/black/{1.5,6}/90, 5/black/{3,6}/90, 6/black/{4.5,6}/90,
    7/black/{6,6}/45, 8/black/{6,4.5}/0, 9/white/{6,3}/0,
    10/black/{1.5,4.5}/10, 11/black/{4.5,4.5}/170}%
  \dzg@edges{3}{(2) -- (3) -- (4) -- (5) -- (6) -- (7) -- (8) (3) -- (10) (7) -- (11)}%
  \dzg@edges{4}{(1) -- (2) (9) -- (8)}}

\@namedef{dzg@sterk@5}{% the barbell roots are unlabelled in Sterk V
  \dzg@barbelllabelsfalse
  \dzg@squareperimeter{4}{2}{1}{nw}{black}{black}%
  \dzg@barbell{7/12, 5/11, 3/10, 1/9}{13}{14}{3}{0.5}}

% -----------------------------------------------------------------------------
% Vinberg diagrams of the 2-elementary lattices (18,2,0) and (18,0,0).
% Doubled and white vertices as in Nikulin's convention for delta.
% -----------------------------------------------------------------------------
% `root labels=ell` labels the twenty boundary lines l_1..l_20 of the integral-
% affine sphere; the two barbell roots carry no line and stay unlabelled.
\tikzset{pics/vinberg 18 2 0/.style={/tikz/transform shape, code={%
  \tikzset{transform shape=false}%
  \ifx\dzg@labelscheme\dzg@ellscheme\dzg@barbelllabelsfalse\fi
  \dzg@squareperimeter{4}{4}{1}{nw}{doubled}{white}%
  \dzg@barbell{13/20, 9/19, 5/18, 1/17}{22}{21}{3}{0.5}}}}

\tikzset{pics/vinberg 18 0 0/.style={/tikz/transform shape, code={%
  \tikzset{transform shape=false}%
  \dzg@vertices{1/doubled/{2,0}/270, 2/white/{1,0}/270, 3/doubled/{0,0}/270,
    4/white/{0,1}/170, 5/doubled/{0,2}/170, 6/white/{0,3}/170,
    7/doubled/{0,4}/170, 8/white/{0,5}/170, 9/doubled/{0,6}/100,
    10/white/{1,5}/10, 11/doubled/{2,4}/10, 12/white/{3,3}/10,
    13/doubled/{4,2}/10, 14/white/{5,1}/10, 15/doubled/{6,0}/270,
    16/white/{5,0}/270, 17/doubled/{4,0}/270, 18/white/{1,1}/10, 19/white/{4,1}/170}%
  \dzg@chain{3}{1}{17}%
  \dzg@edges{3}{(3) -- (18) (15) -- (19)}}}}

% -----------------------------------------------------------------------------
% Coxeter diagrams for the moduli F_2 and F_4 of degree-2 and degree-4 K3
% surfaces.
% -----------------------------------------------------------------------------
\tikzset{pics/k3 F2/.style={/tikz/transform shape, code={%
  \tikzset{transform shape=false}%
  \dzg@trianglepoints{0,0}{7,0}{3.5,6}{6}{6}{6}%
  \foreach \dzg@k in {0,...,5}{%
    \pgfmathtruncatemacro\dzg@j{\dzg@k+6}%
    \pgfmathtruncatemacro\dzg@l{\dzg@k+12}%
    \ifnum\dzg@k=0
      \dzg@vertex{0}{black}{ab-0}{below left}%
      \dzg@vertex{6}{black}{bc-0}{below right}%
      \dzg@vertex{12}{black}{ca-0}{above}%
    \else
      \dzg@vertex{\dzg@k}{black}{ab-\dzg@k}{below}%
      \dzg@vertex{\dzg@j}{black}{bc-\dzg@k}{above right}%
      \dzg@vertex{\dzg@l}{black}{ca-\dzg@k}{above left}%
    \fi}%
  \dzg@vertex{18}{black}{$(A)!2.5/10!(9)$}{below right}%
  \dzg@vertex{21}{black}{$(A)!4.5/10!(9)$}{below right}%
  \dzg@vertex{19}{black}{$(B)!2.5/10!(15)$}{below left}%
  \dzg@vertex{22}{black}{$(B)!4.5/10!(15)$}{below left}%
  \dzg@vertex{20}{black}{$(C)!2.5/10!(3)$}{left}%
  \dzg@vertex{23}{black}{$(C)!4.5/10!(3)$}{left}%
  \dzg@chain{3}{0}{17}%
  \dzg@edges{3}{(17) -- (0) (0) -- (18) (6) -- (19) (12) -- (20)}%
  \dzg@edges{infinity}{
    (18) -- node[vertex label, midway, above, sloped] {$\infty$} (21)
    (21) -- node[vertex label, midway, above, sloped, near end] {$\infty$} (9)
    (19) -- node[vertex label, midway, above, sloped] {$\infty$} (22)
    (22) -- node[vertex label, midway, above, sloped, near end] {$\infty$} (15)
    (20) -- node[vertex label, midway, right] {$\infty$} (23)
    (23) -- node[vertex label, midway, right, near end] {$\infty$} (3)}%
  \dzg@edges{dotted}{(21) -- (22) -- (23) -- (21)}}}}

\tikzset{pics/k3 F4/.style={/tikz/transform shape, code={%
  \tikzset{transform shape=false}%
  \dzg@squareperimeter{6}{4}{0}{sw}{black}{black}%
  \dzg@vertex{16}{black}{$(0)!1/8!(8)$}{below right}%
  \dzg@vertex{17}{black}{$(4)!1/8!(12)$}{below left}%
  \dzg@vertex{18}{black}{$(8)!1/8!(0)$}{above left}%
  \dzg@vertex{19}{black}{$(12)!1/8!(4)$}{above right}%
  \coordinate (10-11) at ($(10)!0.5!(11)$);
  \coordinate (1-2) at ($(1)!0.5!(2)$);
  \coordinate (9-10) at ($(9)!0.5!(10)$);
  \coordinate (2-3) at ($(2)!0.5!(3)$);
  \dzg@vertex{23}{black}{$(10-11)!1/3!(1-2)$}{above left}%
  \dzg@vertex{20}{black}{$(10-11)!2/3!(1-2)$}{below left}%
  \dzg@vertex{22}{black}{$(9-10)!1/3!(2-3)$}{above right}%
  \dzg@vertex{21}{black}{$(9-10)!2/3!(2-3)$}{below right}%
  \dzg@edges{3}{(0) -- (16) (4) -- (17) (8) -- (18) (12) -- (19)
    (16) -- (23) (16) -- (21) (17) -- (20) (17) -- (22)
    (18) -- (23) (18) -- (21) (19) -- (22) (19) -- (20)}%
  \dzg@edges{infinity}{
    (23) -- node[vertex label, midway, right] {$\infty$} (20)
    (22) -- node[vertex label, midway, left] {$\infty$} (21)}}}}

% -----------------------------------------------------------------------------
% Finite and affine Coxeter diagrams drawn with multiple edges: an edge of
% weight m is drawn with m-2 lines.
% -----------------------------------------------------------------------------
\newcommand{\coxeterHthree}{%
  \begin{dynkinDiagram}[label, arrows=false, edge length=.75cm]A{***}
    \dynkinTripleEdge 12
  \end{dynkinDiagram}}
\newcommand{\coxeterHfour}{%
  \begin{dynkinDiagram}[label, arrows=false, edge length=.75cm]A{****}
    \dynkinTripleEdge 12
  \end{dynkinDiagram}}
\newcommand{\coxeterGtwo}{%
  \begin{dynkinDiagram}[label, arrows=false, edge length=1.75cm]A{**}
    \dynkinQuadrupleEdge 12
  \end{dynkinDiagram}}
\newcommand{\coxeterGtwoAffine}{%
  \begin{dynkinDiagram}[label, arrows=false, edge length=1.75cm, labels={0,1,2}]A{***}
    \dynkinQuadrupleEdge 23
  \end{dynkinDiagram}}

% -----------------------------------------------------------------------------
% Table legends: two roots and the edge between them, or one vertex.
% -----------------------------------------------------------------------------
%   \coxeterpair[marks=black/white, edge=4, roots=f_1/f_2, weight=m]
% marks: the two vertex marks (default white/white); edge: a coxeter edge
% value or none (default 3); roots: the two labels in math mode (default
% H_1/H_2); weight: text set above the edge (default empty).
%   \coxetervertex{white|black|doubled}
\pgfqkeys{/dzg/coxeter pair}{
  marks/.store in=\dzg@pairmarks, marks=white/white,
  edge/.store in=\dzg@pairedge, edge=3,
  roots/.store in=\dzg@pairroots, roots=H_1/H_2,
  weight/.store in=\dzg@pairweight, weight=,
}
\def\dzg@none{none}
\def\dzg@splitpair#1/#2\@nil#3#4{\def#3{#1}\def#4{#2}}
\newcommand{\coxeterpair}[1][]{%
  \begin{tikzpicture}[baseline=-0.5ex]
    \pgfqkeys{/dzg/coxeter pair}{#1}%
    \expandafter\dzg@splitpair\dzg@pairmarks\@nil\dzg@marka\dzg@markb
    \expandafter\dzg@splitpair\dzg@pairroots\@nil\dzg@roota\dzg@rootb
    \node[vertex=\dzg@marka, label={[vertex label]above:$\dzg@roota$}] (h1) at (0,0) {};
    \node[vertex=\dzg@markb, label={[vertex label]above:$\dzg@rootb$}] (h2) at (3,0) {};
    \ifx\dzg@pairedge\dzg@none\else
      \dzg@edges{\dzg@pairedge}{(h1) -- node[vertex label, midway, above] {\dzg@pairweight} (h2)}%
    \fi
  \end{tikzpicture}}
\newcommand{\coxetervertex}[1]{\tikz[baseline=-0.5ex]\node[vertex=#1] {};}


